\usepackage{tikz}
\usepackage{caption}
\setbeamertemplate{itemize item}{\textbullet} % Cambiar el símbolo de ítem a punto
\setbeamertemplate{itemize subitem}{\textbullet} % Cambiar el símbolo de subítem a triángulo negro
\setbeamertemplate{itemize subitem}{\blacktriangleright}
\setbeamercolor{frametitle}{bg=blue!20, fg=black}  % Color del título de las diapositivas
% Definición del bloque
\newtcolorbox{myblock}[2][]{colback=blue!80!black!9, colframe=blue!100!black!15,
    fonttitle=\color{black}\bfseries, title=#2, #1, rounded corners=all, boxrule=0.8mm}
\setbeamercolor{frametitle}{bg=blue!20, fg=black}  % Color del título de las diapositivas
\usepackage{fancybox}
\usepackage{tcolorbox}
\usepackage{xcolor}
\usepackage{fancyhdr}

## Ejemplo de Lista

\begin{itemize}
    \item[$\blacktriangleright$] Tercera identidad macroeconómica:
    \item[$\blacktriangleright$] Cuarta identidad macroeconómica:
    \item[$\blacktriangleright$] Quinta identidad macroeconómica:
\end{itemize}


% Definición del bloque
\newtcolorbox{myblock}[2][]{colback=blue!4, colframe=blue!15,
    fonttitle=\color{black}\bfseries, title=#2, #1, rounded corners=all, boxrule=0.8mm}
\setbeamercolor{frametitle}{bg=blue!15, fg=black}  % Color del título de las diapositivas

% Comando para subtítulos sin recuadro
\newcommand{\mySubtitle}[1]{%
  \vspace{0.5cm} % Espacio antes del subtítulo
  \textnormal{\large \textbf{#1}} % Ajustar el formato que desees
  \vspace{0.5cm} % Espacio después del subtítulo
}
